\chapter{L'équipe d'apprentissage automatique}\label{ch:ml:the-machine-learning-team}

Un modèle qui avait demandé trois semaines à un chercheur a mis cinq mois à entrer en trading. Rien, dans ce modèle, n'était difficile. Chaque transmission entre les
personnes qui l'avaient construit formait une file d'attente~: le comité de revue se réunissait chaque semaine, les ingénieurs qui avaient réécrit les variables pour la production accumulaient un arriéré, les
validateurs le renvoyaient deux fois parce que la réécriture ne correspondait pas à la recherche, et le trading simulé (paper trading) ne pouvait pas être
accéléré. Ce chapitre porte sur l'organisation qui entoure les modèles des vingt-sept chapitres précédents~: qui fait quoi, ce qui
circule entre eux, et comment mesurer le résultat. Il construit deux objets. Le premier est une \omterm{def:ml:the-machine-learning-team:handoff}{spécification de transmission} que peut vérifier une
tâche d'intégration continue, de sorte que le paquet remis par un chercheur s'exécute en production comme il s'exécutait en recherche, ou
échoue à un test nommé. Le second est un simulateur du pipeline recherche-production, modélisé comme un réseau de files d'attente. Sur le pipeline du
chapitre, le modèle médian atteint la production 20 semaines après la fin de la recherche, et 65~\% des modèles sont perdus en chemin, pour la plupart à la
première revue~; remplacer la réécriture par un paquet validé ramène la médiane à 11 semaines, dont une semaine correspond au temps que le chercheur passe
à construire le paquet.

\section{Rôles~: chercheur, ingénieur, plateforme}

\begin{definition}[Ingénieur en apprentissage automatique, propriétaire du modèle]\label{def:ml:the-machine-learning-team:roles}
Un \emph{ingénieur en apprentissage automatique}\index{ingénieur en apprentissage automatique} (ML engineer) conçoit et exploite les systèmes qui entraînent, servent et surveillent
les modèles en production, en combinant les compétences des data scientists, des data engineers et des ingénieurs logiciels. Le \emph{propriétaire du
modèle}\index{propriétaire du modèle} est la personne nommément responsable d'un modèle en production~: sa performance, sa surveillance, son
réentraînement et son retrait, ainsi que les réponses aux questions que lui posent les validateurs, les responsables des risques et les régulateurs.
\end{definition}

Le travail d'apprentissage automatique d'une société de trading mobilise au moins trois catégories de personnes. Les chercheurs (chercheurs quantitatifs ou
data scientists) identifient le signal et ajustent le modèle, avec les méthodes de ce livre et le savoir-faire du livre~7. Les \omterm{def:ml:the-machine-learning-team:roles}{ingénieurs en apprentissage
automatique} transforment le modèle en un service qui fonctionne chaque jour sans son auteur~: variables calculées en ligne comme elles l'étaient hors ligne (chapitre~24),
modèle compilé dans le budget de latence (chapitre~26), dispositifs de surveillance (chapitre~27). Une équipe plateforme construit ce que tous les modèles partagent. Kreuzberger, Kühl et Hirschl,
à partir d'une revue de la littérature, d'une revue d'outils et d'entretiens avec des experts, recensent sept rôles autour des opérations d'apprentissage automatique (partie prenante métier,
architecte de solutions, data scientist, data engineer, ingénieur logiciel, ingénieur DevOps, et l'\omterm{def:ml:the-machine-learning-team:roles}{ingénieur en apprentissage automatique} ou MLOps qui combine plusieurs de ces rôles)~;
une société de trading ajoute le validateur de modèle du livre~6 (chapitre~26), indépendant des autres, ainsi que le trader ou le gérant de portefeuille
qui utilise le modèle.

Le \omterm{def:ml:the-machine-learning-team:roles}{propriétaire du modèle} est un rôle, non une équipe. Il s'agit en général du chercheur qui a proposé le modèle ou du responsable du desk qui l'utilise~; ce qui compte
c'est qu'il s'agisse d'une personne nommée, inscrite dans le \omterm{def:ml:experiment-tracking-and-reproducibility:registry}{registre de modèles} (chapitre~25) et sur la \omterm{def:ml:interpretability-and-model-governance:card}{fiche de modèle} (chapitre~21), qui répond à l'alerte d'une
\omterm{def:ml:monitoring-and-retraining:performance-alarm}{alarme de performance} et signe la décision de réentraîner ou de retirer le modèle. Un modèle sans propriétaire est celui dont les dispositifs de surveillance n'alertent personne. Amershi et ses coauteurs, en observant des équipes de Microsoft, ont constaté
que la personnalisation et la réutilisation des modèles demandent des compétences différentes de celles des équipes logicielles, et que les modèles se prêtent moins bien à un traitement en modules distincts
que les composants logiciels ordinaires~; ces deux constats expliquent pourquoi la répartition du travail entre chercheurs et ingénieurs est un choix de conception, avec un coût,
et non une donnée de départ.

\section{Comment le travail circule~: la transmission}

\begin{definition}[Spécification de transmission]\label{def:ml:the-machine-learning-team:handoff}
Une \emph{spécification de transmission}\index{spécification de transmission} (handoff) est le contrat entre les personnes qui construisent un modèle et celles qui le font
fonctionner~: la liste de ce que doit contenir un paquet de modèle (l'artefact et son empreinte de hachage, les \omterm{def:ml:data-and-feature-stores:store}{définitions de variables}, des vecteurs de test
d'entrées avec les sorties que le modèle doit reproduire, le budget de latence accompagné d'une mesure, la spécification de surveillance,
la \omterm{def:ml:interpretability-and-model-governance:card}{fiche de modèle}) et les contrôles que le paquet doit passer avant que la production ne l'accepte.
\end{definition}

Dans la société de l'accroche, la transmission était une réécriture. Le chercheur remettait un notebook~; un ingénieur réimplémentait les variables et le modèle
en code de production~; le validateur comparait les deux et, en cas d'écart, renvoyait le modèle. Chaque écart représentait un aller-retour à travers deux files d'attente. Les règles de Zinkevich pour l'ingénierie de l'apprentissage
automatique le disent sans détour~: «~Réutilisez le code entre votre pipeline d'entraînement et votre pipeline de service chaque fois que c'est possible~». La \omterm{def:ml:the-machine-learning-team:handoff}{spécification de transmission} est l'autre moitié de ce conseil~: lorsque le même code
tourne aux deux endroits, ce qui circule entre les équipes est
un paquet qui peut être vérifié mécaniquement.

Le paquet du chapitre concerne la forêt compilée du chapitre~26 (le \cref{lst:ml:team:validate}). Son manifeste nomme le fichier de la forêt et
son empreinte de hachage (\texttt{firm.workflow} du livre~7)~; seize \omterm{def:ml:data-and-feature-stores:store}{définitions de variables} au format de \texttt{firm.featstore}~; les valeurs de la recherche pour ces variables sur une session de référence
du carnet d'ordres simulé~; 200 vecteurs de test avec les sorties de la forêt, à reproduire à $10^{-12}$ près~; un budget de latence de 5 microsecondes, avec une valeur mesurée de 2.6~;
quatre dispositifs de surveillance, chacun avec des seuils, un taux de fausses alertes et un propriétaire~; et les quinze champs de la \omterm{def:ml:interpretability-and-model-governance:card}{fiche de modèle}.
Le validateur effectue neuf contrôles, et le paquet propre les passe tous. Le chapitre introduit ensuite huit défauts que porte habituellement une transmission,
un à la fois (le \cref{tab:ml:team:defects}).

\begin{table}[htbp]
\centering\small
\begin{tabular}{@{}p{5.4cm}p{6.6cm}@{}}
\toprule
défaut introduit & contrôles qui échouent\\
\midrule
aucun (le paquet tel que remis) & aucun des neuf\\
artefact reconstruit après son hachage & empreinte de l'artefact~; vecteurs de test\\
une variable lit une entrée que le flux ne transporte pas & spécification des variables\\
une fenêtre de 5 secondes écrite 5\,000 (millisecondes) & vecteurs de variables\\
une variable manquante (15 pour un modèle qui en lit 16) & spécification des variables\\
vecteurs de test contenant les sorties d'un autre modèle & vecteurs de test\\
latence de l'appel de bibliothèque, et non du modèle compilé & latence\\
un dispositif de surveillance sans propriétaire & surveillance\\
une \omterm{def:ml:interpretability-and-model-governance:card}{fiche de modèle} sans limites & \omterm{def:ml:interpretability-and-model-governance:card}{fiche de modèle}\\
\bottomrule
\end{tabular}
\caption{Le validateur de transmission sur le paquet propre et sur huit défauts introduits. Données~: \texttt{ml\_team.handoff\_table}.}
\label{tab:ml:team:defects}
\end{table}

Chaque défaut est détecté, et par le contrôle écrit pour lui, avec une leçon dans la troisième ligne. Une fenêtre exprimée dans la mauvaise unité reste une
\omterm{def:ml:data-and-feature-stores:store}{définition de variable} valide, et les calculs hors ligne et en ligne concordent entre eux, puisque les deux utilisent la même fenêtre erronée~: les contrôles de spécification et de parité réussissent.
Seule la comparaison avec les valeurs de variables calculées par la recherche sur une session de référence la révèle.
Les vecteurs de test doivent couvrir les variables autant que le modèle.

\section{Ce que la plateforme possède}

\begin{definition}[Plateforme d'apprentissage automatique]\label{def:ml:the-machine-learning-team:platform}
Une \emph{plateforme d'apprentissage automatique}\index{plateforme d'apprentissage automatique} est l'infrastructure commune sur laquelle chaque modèle est construit
et exécuté~: magasins de données et \omterm{def:ml:data-and-feature-stores:store}{feature stores}, infrastructure d'entraînement, suivi des expériences et \omterm{def:ml:experiment-tracking-and-reproducibility:registry}{registre de modèles}, compilation et service,
surveillance, ainsi que les contrôles qu'un paquet de modèle passe avant d'aller en production, gérée par une seule équipe et utilisée par toutes.
\end{definition}

La plateforme possède tout ce qui est identique pour chaque modèle. Dans les termes de ce livre~: le \omterm{def:ml:data-and-feature-stores:store}{feature store} et ses tests de parité (chapitre~24),
l'infrastructure d'entraînement (chapitre~23), le suivi, le \omterm{def:ml:experiment-tracking-and-reproducibility:registry}{registre de modèles} et la \omterm{def:ml:experiment-tracking-and-reproducibility:registry}{piste d'audit} (chapitre~25), les compilateurs et les noyaux de service (chapitre~26),
les dispositifs de surveillance (chapitre~27) et le validateur de transmission.
Les chercheurs possèdent le modèle et ses variables~; la plateforme possède la garantie qu'une variable a le même sens en recherche et en production. La frontière compte, car chaque côté tend à faire passer du travail de l'autre côté~:
les chercheurs demandent des pipelines sur mesure, et les plateformes accumulent des variables qu'aucun modèle n'utilise. Un bon test est la quatrième règle de Zinkevich, «~Gardez le premier modèle simple et faites bien l'infrastructure~»~:
une plateforme se justifie par le second modèle, non par le premier. Sa seconde justification tient à la visibilité~:
parmi les risques propres aux systèmes d'apprentissage automatique, Sculley et ses coauteurs citent les consommateurs non déclarés
et les dépendances aux données, que rendent explicites un \omterm{def:ml:experiment-tracking-and-reproducibility:registry}{registre de modèles} partagé et un \omterm{def:ml:data-and-feature-stores:store}{feature store}.

Le validateur est la pièce la plus rentable de la plateforme, car il transforme une conversation en test. Dans le pipeline du chapitre décrit ci-dessous,
c'est lui qui permet d'intégrer un modèle en une semaine au lieu d'une réécriture de quatre semaines.

\section{Mesurer le pipeline}

Le pipeline est un réseau de files d'attente (le \cref{lst:ml:team:sim})~; ses paramètres sont illustratifs et ne sont mesurés dans aucune société. Les modèles arrivent à la fin de la recherche, à raison de 0.35 par semaine (18 par an), et
traversent cinq postes. Le comité de revue (une place, 0.6 semaine par modèle) élimine 35~\% des modèles qu'il examine et renvoie 15~\% au chercheur
pour deux semaines de corrections de recherche. Deux ingénieurs réimplémentent chaque modèle retenu en environ quatre semaines~; 30~\% des réécritures reviennent au chercheur
pour une semaine et demie de précisions, puis sont réécrites, et 5~\% sont abandonnées. Un validateur (1.5 semaine)
élimine 10~\% et renvoie 25~\% aux ingénieurs. Le trading simulé dure quatre semaines et élimine 20~\%~; un test en canari de deux semaines élimine 5~\%. Les durées de service
sont exponentielles, à l'exception des deux périodes de trading fixes~; les 500 premiers des 4\,000 modèles simulés sont écartés
pendant que le pipeline se remplit.

Le simulateur est vérifié sur le seul cas à forme fermée. Une station unique de la file M/M/1 du livre~4 (chapitre~8), avec des arrivées
à 0.8 et un service à 1 par semaine, donne un temps moyen dans le système de 5.05 semaines, contre la valeur théorique $1/(\mu-\lambda)=5$~; la
loi stationnaire de la chaîne de naissance et de mort correspondante, calculée avec \texttt{firm.queues} du livre~4, donne un nombre moyen de 4.00 modèles dans
le système et un système vide 20~\% du temps, et le taux d'arrivée simulé multiplié par son temps moyen donne 4.04, comme le veut la loi de Little
$L=\lambda W$ l'exige. Sur l'ensemble du pipeline, le nombre moyen de modèles présents dans celui-ci est de 4.87 et le taux d'arrivée multiplié par leur
temps moyen passé dans celui-ci (14.0 semaines, en comptant les modèles éliminés jusqu'à leur sortie) vaut 4.90.

\begin{table}[htbp]
\centering\small
\begin{tabular}{@{}lcc@{}}
\toprule
 & réécriture (transmission) & paquet validé\\
\midrule
délai médian jusqu'en production (semaines) & 20.4 & 11.1\\
quartiles (semaines) & 14.8 à 30.2 & 9.6 à 13.7\\
90\up{e}~centile (semaines) & 42.9 & 16.5\\
part des modèles atteignant la production & 35.1~\% & 37.5~\%\\
perdus à la revue & 41.1~\% & 39.5~\%\\
perdus à la construction (réimplémentation / intégration) & 5.3~\% & 3.5~\%\\
perdus à la validation & 7.9~\% & 6.9~\%\\
perdus en trading simulé & 8.7~\% & 10.3~\%\\
perdus dans le canari & 1.9~\% & 2.2~\%\\
passages par la construction, par modèle ayant franchi la revue & 1.79 & 1.14\\
taux d'utilisation des ingénieurs & 74~\% & 12~\%\\
\bottomrule
\end{tabular}
\caption{Le pipeline de la recherche à la production avec une réécriture à la transmission et avec un paquet validé (le délai inclut une
semaine de recherche consacrée à la construction du paquet)~; 3\,500 modèles chacun. Données~: \texttt{ml\_team.pipeline}.}\label{tab:ml:team:pipeline}
\end{table}

Le \Cref{tab:ml:team:pipeline} donne le résultat du chapitre. Avec la réécriture, le modèle médian atteint la production 20.4 semaines après
la fin de la recherche, soit environ cinq mois, et un modèle sur dix met plus de 42.9 semaines~; 35~\% des modèles y parviennent. La perte la plus importante se situe à la
revue, 41~\% de l'ensemble des modèles~: c'est là que la plupart des modèles doivent mourir, à moindre coût et au plus tôt. Sur le temps total, le service sans aucune attente
représenterait 14.9 semaines du délai médian~; le reste est dû à la file d'attente, concentrée chez les ingénieurs, occupés 74~\% du temps
car chaque modèle passe par eux 1.79 fois en moyenne. Remplacer la réécriture par un paquet validé (intégration en une semaine environ,
reprises à la construction ramenées de 30~\% à 5~\% et à la validation de 25~\% à 10~\%, plus une semaine du temps du chercheur)
ramène le délai médian à 11.1 semaines, soit un gain de 9.3 semaines (45~\%), et le 90\up{e}~centile de 42.9 à 16.5. Les parts perdues à chaque
porte bougent à peine, ce qui est normal~: supprimer une transmission doit rendre les décisions plus rapides, non pas différentes.

\begin{omfigure}
\begin{tikzpicture}
\begin{axis}[width=11cm,height=5.8cm,xlabel={modèles arrivant par semaine},ylabel={délai médian (semaines)},xmin=0.19,xmax=0.45,
  ymin=0,ymax=72,tick label style={font=\small},label style={font=\small},grid=major,grid style={gray!25},legend pos=north west,
  legend cell align=left,legend style={font=\small},xticklabel style={/pgf/number format/fixed,/pgf/number format/precision=2},
  table/col sep=comma]
\addplot[omProp,very thick,mark=*,mark size=1.5pt] table[x=rate,y=handoff]{figdata/ml/28-the-machine-learning-team/rate.csv};
\addplot[omDef,very thick,dashed,mark=square*,mark size=1.5pt,mark options={solid}] table[x=rate,y=package]{figdata/ml/28-the-machine-learning-team/rate.csv};
\legend{réécriture à la transmission,paquet validé}
\end{axis}
\end{tikzpicture}

\omcaption{Délai médian jusqu'en production en fonction du rythme auquel la recherche envoie des modèles, avec une réécriture à la transmission et
avec un paquet validé, moyennes sur trois pipelines simulés de 4\,000 modèles. Données~: \texttt{ml\_team.rate\_curve}.}
\label{fig:ml:team:rate}
\end{omfigure}

Le tableau tient pour un seul rythme d'arrivée~; la \cref{fig:ml:team:rate} montre pourquoi cela compte plus qu'il n'y paraît. La station des ingénieurs est
une file proche de la saturation, et le temps d'attente d'une file croît sans borne lorsque son taux d'utilisation s'approche de l'unité~: de 0.35 à 0.44
modèles par semaine, soit un quart de recherche en plus, le délai médian avec la réécriture passe d'environ 22 à 70 semaines, alors qu'avec le
paquet il reste voisin de 11. Un pipeline se mesure par son délai, par la part perdue à chaque porte et par le taux d'utilisation de sa
station la plus chargée, et un responsable qui ajoute des chercheurs sans regarder ce dernier chiffre rend chaque modèle plus lent.

\begin{method}[Faire tourner un pipeline de la recherche à la production]\label{met:ml:team:method}
\begin{enumerate}
\item Désigner, pour chaque modèle, un \omterm{def:ml:the-machine-learning-team:roles}{propriétaire du modèle}, inscrit dans le registre et sur la \omterm{def:ml:interpretability-and-model-governance:card}{fiche de modèle}.
\item Remplacer les réécritures par des paquets~: le même code de variables et de modèle en recherche et en production, et une \omterm{def:ml:the-machine-learning-team:handoff}{spécification de transmission} (handoff)
vérifiée en intégration continue, avec des vecteurs de test pour les variables comme pour le modèle.
\item Éliminer les modèles tôt~: la porte la moins coûteuse doit rejeter le plus.
\item Mesurer le délai, les pertes à chaque porte et le taux d'utilisation de chaque équipe~; garder la station la plus chargée bien en deçà de la saturation.
\item Mettre ce que partagent tous les modèles sur une plateforme possédée par une seule équipe, et la justifier dès le second modèle.
\end{enumerate}
\end{method}

\section{Tutoriel~: qui possède le modèle~?}\label{tut:ml:the-machine-learning-team:tut}

\begin{tutorial}
\textbf{Objectif.} Valider un paquet de transmission comme le ferait l'intégration continue, introduire des défauts et voir quel contrôle détecte chacun~;
simuler le pipeline avec et sans réécriture. \textbf{État final~:} \cref{tab:ml:team:defects}, \cref{tab:ml:team:pipeline},
\cref{fig:ml:team:rate}.
\begin{tutsteps}
\item \textbf{Le validateur.}
\omcode{code/firm/modelpkg/firm_modelpkg.py}{60}{92}{Schéma, empreinte de l'artefact, spécification des variables, parité et vecteurs de variables.}\label{lst:ml:team:validate}
\item \textbf{Les autres contrôles.}
\omcode{code/firm/modelpkg/firm_modelpkg.py}{93}{105}{Vecteurs de test passés dans le noyau de forêt, latence, surveillance et fiche de modèle.}\label{lst:ml:team:checks}
\item \textbf{Les défauts introduits} (\texttt{ml\_team.\_defects}).
\omcode{code/ml/28-the-machine-learning-team/python/ml_team.py}{96}{119}{Sept des huit défauts, chacun étant une modification d'un paquet chargé.}\label{lst:ml:team:defects}
\item \textbf{Le pipeline.}
\omcode{code/ml/28-the-machine-learning-team/python/ml_team.py}{145}{154}{Stations, avec et sans réécriture.}\label{lst:ml:team:stations}
\item \textbf{Le simulateur.}
\omcode{code/firm/modelpkg/firm_modelpkg.py}{143}{171}{La boucle d'événements~: files, service, et les décisions d'élimination, de reprise et de passage.}\label{lst:ml:team:sim}
\item \textbf{Exécuter} \texttt{ml\_team.handoff\_table()}, \texttt{pipeline(True)}, \texttt{pipeline(False, 1.0)},
\texttt{mm1\_check()} et \texttt{fig\_team.py} (environ 20 secondes sur un cœur).
\end{tutsteps}
\textbf{Pour aller plus loin.} Faire se réunir le comité de revue deux fois plus souvent, puis trouver quelle station devient alors la plus chargée.
\end{tutorial}

\section{Construction~: le paquet de modèle}\label{bld:ml:the-machine-learning-team:modelpkg}

\begin{build}
\bfield{Rôle} Un contrat entre la recherche et la production qu'une machine peut vérifier, et un modèle du pipeline qu'il traverse.
\bfield{Interface} \texttt{REQUIRED}, \texttt{INPUTS}, \texttt{MONITORS}~; \texttt{load(path)}~;
\texttt{validate(pkg, events, times)} renvoyant des triplets (contrôle, réussite, détail), sur \texttt{firm.workflow}, \texttt{firm.featstore}, \texttt{firm.mlinfer} et
\texttt{firm.modelcard}~; \texttt{Station(name, servers, mean, kill, rework, rework\_to, dist, next\_to)}~;
\texttt{simulate\_pipeline(stations, rate, n, seed)}.
\bfield{Règles} Chaque champ de la spécification est obligatoire~; un contrôle qui ne peut pas s'exécuter échoue~; les valeurs des variables, et pas seulement les sorties du modèle,
sont comparées à celles de la recherche~; le simulateur est déterministe étant donné sa graine.
\bfield{Tests d'acceptation} \texttt{code/firm/modelpkg/tests/}~: un paquet complet réussit~; un champ manquant, une mauvaise empreinte,
une entrée inconnue et un vecteur de test périmé échouent chacun à leur contrôle~; une station M/M/1 unique concorde avec $1/(\mu-\lambda)$~; une ligne à retard pur
renvoie son temps fixe~; les probabilités d'élimination et de reprise sont respectées.
\bfield{Pour aller plus loin} Exécuter le validateur comme tâche de pré-fusion sur le registre, avec le manifeste du paquet comme artefact enregistré~;
fixer des priorités à la station la plus chargée.
\end{build}

\begin{omsources}
\item S.~Amershi et coauteurs, ``Software engineering for machine learning: a case study'', \emph{ICSE Software Engineering
in Practice}, 2019.
\item D.~Kreuzberger, N.~Kühl et S.~Hirschl, ``Machine learning operations (MLOps): overview, definition, and
architecture'', \emph{IEEE Access}, 2023.
\item M.~Zinkevich, \emph{Rules of Machine Learning: Best Practices for ML Engineering}, Google for Developers.
\item J.~D.~C.~Little, ``A proof for the queuing formula $L=\lambda W$'', \emph{Operations Research}, 1961.
\item D.~Sculley et coauteurs, ``Hidden technical debt in machine learning systems'', \emph{NeurIPS}, 2015.
\end{omsources}

\section{Exercices}

\begin{exercise}[$\star$]\label{exo:ml:the-machine-learning-team:1}
Nommer les trois catégories de personnes intervenant dans le travail d'apprentissage automatique d'une société de trading et une chose que chacune possède.
\end{exercise}

\begin{exercise}[$\star$]\label{exo:ml:the-machine-learning-team:2}
Pourquoi le \omterm{def:ml:the-machine-learning-team:roles}{propriétaire du modèle} doit-il être une personne nommément désignée~? Que se passe-t-il lorsqu'il s'agit d'une équipe~?
\end{exercise}

\begin{exercise}[$\star$]\label{exo:ml:the-machine-learning-team:3}
Vérifier la loi de Little sur la station unique du chapitre~: des arrivées à 0.8 par semaine, un temps moyen dans le système de 5.05 semaines. Quel est le
nombre moyen de modèles dans le système~?
\end{exercise}

\begin{exercise}[$\star\star$]\label{exo:ml:the-machine-learning-team:4}
Pourquoi les contrôles de spécification et de parité laissent-ils passer une fenêtre écrite en millisecondes, et quel contrôle le détecte~?
\end{exercise}

\begin{exercise}[$\star\star$]\label{exo:ml:the-machine-learning-team:5}
Les ingénieurs sont occupés 74~\% du temps à la réécriture. En utilisant la formule de M/M/1 comme approximation grossière, de combien le temps
d'attente moyen change-t-il si leur taux d'utilisation passe à 90~\%~?
\end{exercise}

\begin{exercise}[$\star\star$]\label{exo:ml:the-machine-learning-team:6}
\emph{Trouver la faille.} «~Notre pipeline est lent parce que la validation élimine trop de modèles~; il faut assouplir les validateurs.~»
\end{exercise}

\begin{exercise}[$\star\star\star$]\label{exo:ml:the-machine-learning-team:7}
\emph{Programmation.} Conserver la réécriture et donner un troisième membre à l'équipe des ingénieurs. Comparer le délai de mise en production médian avec celui du paquet, et
indiquer la modification à apporter en premier.
\end{exercise}

\begin{exercise}[$\star\star\star$]\label{exo:ml:the-machine-learning-team:8}
Rédiger la \omterm{def:ml:the-machine-learning-team:handoff}{spécification de transmission} (handoff) du modèle linéaire surveillé du chapitre~27~: les champs, les vecteurs de test et les contrôles.
\end{exercise}

\section{Problème~: Qui possède le modèle~?}

\begin{problem}[{Problème du week-end --- qui possède le modèle~?}]\label{pb:ml:the-machine-learning-team:1}
Le paquet de transmission et le pipeline du chapitre.

\medskip\noindent\textbf{Partie I --- Rôles.}
\begin{enumerate}
\item Nommer les rôles dans le travail d'apprentissage automatique d'une société de trading.
\item À quoi le \omterm{def:ml:the-machine-learning-team:roles}{propriétaire du modèle} répond-il~?
\item Que possède la plateforme, et que ne possède-t-elle pas~?
\item Pourquoi la répartition entre chercheurs et ingénieurs est-elle un choix de conception~?
\end{enumerate}

\medskip\noindent\textbf{Partie II --- La transmission.}
\begin{enumerate}[resume]
\item Que contient le paquet du chapitre~?
\item Quels contrôles le validateur exécute-t-il~?
\item Quel contrôle détecte chaque défaut introduit~?
\item Pourquoi les vecteurs de test doivent-ils couvrir les variables~?
\end{enumerate}

\medskip\noindent\textbf{Partie III --- Le pipeline.}
\begin{enumerate}[resume]
\item Décrire le pipeline simulé.
\item Comment le simulateur est-il vérifié~?
\item Où meurent les modèles, et où part le temps~?
\item Que se passe-t-il lorsque davantage de modèles arrivent~?
\end{enumerate}

\medskip\noindent\textbf{Partie IV --- Le verdict.}
\begin{enumerate}[resume]
\item Énoncer le \emph{résultat nommé}~: le délai médian de mise en production du pipeline, la part de modèles perdue à chaque contrôle et le gain obtenu en
supprimant une transmission.
\item Pourquoi les pertes à chaque contrôle devraient-elles à peine changer lorsque la transmission est supprimée~?
\item Que paie le chercheur pour le paquet, et cela en vaut-il la peine~?
\item Un troisième ingénieur ferait-il aussi bien~?
\item Que devrait mesurer la société chaque mois~?
\item Qui devrait être le propriétaire de la forêt du chapitre~?
\item Que doit-il se passer lorsqu'un paquet échoue à un contrôle~?
\item En une phrase~: qu'est-ce qui rend une transmission peu coûteuse~?
\end{enumerate}
\end{problem}

\section{Questions d'entretien}

\begin{interviewq}[$\star$ \iqroles{ingénieur ML}]\label{iq:ml:the-machine-learning-team:1}
Que fait un \omterm{def:ml:the-machine-learning-team:roles}{ingénieur en apprentissage automatique} qu'un chercheur ne fait pas~?
\end{interviewq}

\begin{interviewq}[$\star\star$ \iqroles{ingénieur ML, développeur}]\label{iq:ml:the-machine-learning-team:2}
Que doit remettre une équipe de recherche à la production avec un modèle~?
\end{interviewq}

\begin{interviewq}[$\star\star$ \iqroles{ingénieur ML}]\label{iq:ml:the-machine-learning-team:3}
Vos variables diffèrent entre la recherche et la production. Comment le découvrez-vous, et comment empêchez-vous que cela se reproduise~?
\end{interviewq}

\begin{interviewq}[$\star\star$ \iqroles{chercheur, ingénieur ML}]\label{iq:ml:the-machine-learning-team:4}
Les modèles mettent cinq mois à atteindre la production. Comment en trouveriez-vous la raison~?
\end{interviewq}

\begin{interviewq}[$\star\star$ \iqroles{ingénieur ML}]\label{iq:ml:the-machine-learning-team:5}
Qu'est-ce qui doit figurer sur une \omterm{def:ml:the-machine-learning-team:platform}{plateforme d'apprentissage automatique}, et qu'est-ce qui doit rester dans l'équipe du modèle~?
\end{interviewq}

\begin{interviewq}[$\star\star\star$ \iqroles{ingénieur ML}]\label{iq:ml:the-machine-learning-team:6}
Votre équipe plateforme doit doubler le nombre de modèles qu'elle prend en charge l'an prochain. Que mesurez-vous en premier~?
\end{interviewq}
